\section*{अध्याय \ref{ch:b3:orogenic-wedges} --- अधिक्षेप पट्टियाँ और पर्वतजनक कीलें}

\begin{solution}{exo:b3:orogenic-wedges:1}
$\sin\ang{30} = 0.5$: $K = 2\times0.5/0.5 = 2$ और $\sigma'_1/\sigma'_3 = 1 +
K = 3$.
\end{solution}

\begin{solution}{exo:b3:orogenic-wedges:2}
$\phi = \arctan0.6 = \ang{31.0}$, $\sin\phi = 0.514$, $K = 2.12$।
क्षैतिज आधार पर $\alpha = 0.3/3.12 = \qty{0.0962}{rad} = \ang{5.5}$।
नति $\ang{3} = \qty{0.0524}{rad}$ वाले आधार पर, $\alpha + \beta = 0.352/3.12 =
\qty{0.113}{rad} = \ang{6.5}$, इसलिए $\alpha = \ang{3.5}$।
\end{solution}

\begin{solution}{exo:b3:orogenic-wedges:3}
गहराई $= (T - 10)/30$: ऐपेटाइट (U--Th)/He 1.3 से \qty{2.3}{km}; ऐपेटाइट विखंडन-पथ
2.7 से \qty{3.7}{km}; ज़िरकॉन विखंडन-पथ $270/30 = \qty{9.0}{km}$।
\qty{5}{km} पर चट्टान का तापमान $10 + 150 = \qty{160}{\celsius}$ था: ज़िरकॉन
(U--Th)/He, जिसका अवरोध-तापमान 160 से \qty{200}{\celsius} है, 5.0 से
\qty{6.3}{km} गहराई पर।
\end{solution}

\begin{solution}{exo:b3:orogenic-wedges:4}
क्षैतिज लंबाई $2/\tan\ang{30} = \qty{3.46}{km}$; नति के अनुदिश
$2/\sin\ang{30} = \qty{4.0}{km}$। संस्तर \omterm{def:b3:orogenic-wedges:thrust-geometry}{रैंप} की ऊँचाई से उठ जाता है,
\qty{2}{km}। कट-ऑफ़ को पूरे \omterm{def:b3:orogenic-wedges:thrust-geometry}{रैंप} को पार करना होगा, यानी
\qty{4.0}{km} का विस्थापन।
\end{solution}

\begin{solution}{exo:b3:orogenic-wedges:5}
$\phi = \arctan1.03 = \ang{45.8}$, $K = 5.08$; $\beta = \qty{0.105}{rad}$।
$\lambda = 0.4$ के लिए: $\alpha + \beta = (0.105 + 0.85\times0.6)/(1 +
5.08\times0.6) = 0.615/4.05 = \qty{0.152}{rad} = \ang{8.7}$, $\alpha =
\ang{2.7}$। $\lambda = 0.9$ के लिए: $0.190/1.51 = \qty{0.126}{rad} = \ang{7.2}$,
$\alpha = \ang{1.2}$। द्रव दाब बढ़ाने से \omterm{def:b3:orogenic-wedges:critical-taper}{क्रांतिक टेपर} घटता है:
कील अतिक्रांतिक हो जाती है और अपने आधार पर बिना मोटाई बढ़े आगे खिसकती है,
और उसका अग्र भाग तब तक आगे बढ़ता है जब तक \omterm{def:b3:orogenic-wedges:accretion}{अग्र संचयन} और अपरदन उसे
नए टेपर पर वापस ले आते हैं।
\end{solution}

\begin{solution}{exo:b3:orogenic-wedges:6}
$S = A/r = 6.0/1.5 = \qty{4.0}{km}$। \omterm{def:b3:orogenic-wedges:thrust-folds}{विलगन वलन} के रूप में पढ़ने पर, $S = A/h =
6.0/3.0 = \qty{2.0}{km}$: सही मान का आधा। दोनों व्याख्याओं को
वलन के आकार और \omterm{def:b3:orogenic-wedges:thrust-geometry}{रैंप} की उपस्थिति से अलग पहचाना जाना चाहिए।
\end{solution}

\begin{solution}{exo:b3:orogenic-wedges:7}
$K = 2.12$ और $1 + K(1 - 0.4) = 2.27$। समीकरण देता है $\mu_b(1 -
\lambda_b) = (\alpha + \beta)\times2.27 - \beta$, जहाँ $\beta =
\qty{0.0349}{rad}$। जूरा: $0.0611\times2.27 - 0.0349 = 0.10$। दक्षिण:
$0.218\times2.27 - 0.0349 = 0.46$। द्रवस्थैतिक दाब पर बायरली नियम:
$0.85\times0.6 = 0.51$। दक्षिण का आधार बायरली नियम के निकट है;
जूरा का आधार लगभग पाँच गुना कमज़ोर है, जैसा हैलाइट में एक डेकोलमां से
अपेक्षित है, जिसका घर्षण कम होता है, जब द्रव अति-दाबित नहीं होते।
\end{solution}

\begin{solution}{exo:b3:orogenic-wedges:8}
$\kappa = \qty{e-6}{m^2/s} = \qty{31.6}{km^2/Myr}$, और $v =
\qty{1}{km/Myr}$: $\mathrm{Pe} = 30/31.6 = 0.951$। गुणक $0.951/(1 -
e^{-0.951}) = 1.55$ प्रवणता को 30 से \qty{46.5}{\celsius/km} तक बढ़ा देता है।
\qty{65}{\celsius} के लिए: $1 - e^{-\mathrm{Pe}\,z/L} = (55/900)(1 - e^{-0.951})
= 0.0375$, इसलिए $\mathrm{Pe}\,z/L = 0.0382$ और $z = \qty{1.21}{km}$, जबकि
अपरदन के बिना $55/30 = \qty{1.83}{km}$।
\end{solution}

\begin{solution}{exo:b3:orogenic-wedges:9}
$\mathrm dh/\mathrm dt = \qty{0.8}{km}/\qty{1.2}{Myr} = \qty{0.67}{mm/yr}$।
$f = 0.5$ के साथ, $v = 0.67\times(1 - 0.5) = \qty{0.33}{mm/yr}$।
\end{solution}

\begin{solution}{exo:b3:orogenic-wedges:10}
$\alpha = 0$ और $u = 1 - \lambda$ के लिए: $\beta(1 + Ku) = \beta + \mu_b u$,
इसलिए $\beta Ku = \mu_b u$ और $\beta^* = \mu_b/K$, चाहे $u$ कुछ भी हो। ताइवान के लिए,
$0.85/5.08 = \qty{0.167}{rad} = \ang{9.6}$। अधिक तीव्र नति वाले आधार पर क्रांतिक
सतह-प्रवणता ऋणात्मक हो जाती है: समतल या अग्रभूमि की ओर ढलती सतह वाली कोई भी कील
पहले से ही अतिक्रांतिक है, और वह अपने आधार पर बिना मोटाई बढ़े फिसलती है। एक
अधिक्षेप पट्टी वहाँ अपना सामान्य आकार नहीं बढ़ा सकती।
\end{solution}

\begin{solution}{exo:b3:orogenic-wedges:11}
जब $\Lambda \to \infty$ और $v \to 0$ होते हैं, तो $m \to 0$, इसलिए $e^{-mz_c} \to 1$, और
जियोथर्म सीधी रेखा है, $G_0 = G(z_c) = G$: $f \to (G - \Gamma)/G = 1 -
\Gamma/G = 1 - 6.5/30 = 0.78$। तब प्रवणता $v/(1 - f) = vG/\Gamma =
4.6v$। $v = \qty{1}{mm/yr}$ के लिए, $v/2\kappa = \qty{0.0158}{km^{-1}}$।
\qty{5}{km} पर: $m = 0.0158 + \sqrt{0.0158^2 + 1.257^2} = \qty{1.273}{km^{-1}}$,
$e^{-1.273\times1.21} = 0.215$, $f = 40.0\times0.215/44.7 = 0.19$।
\qty{20}{km} पर: $m = 0.0158 + 0.3146 = \qty{0.330}{km^{-1}}$, $e^{-0.399} =
0.671$, $f = 40.0\times0.671/44.7 = 0.60$।
\end{solution}

\begin{solution}{exo:b3:orogenic-wedges:12}
काट एक त्रिभुज है जिसका आधार $W$ और ऊँचाई $W\tan(\alpha + \beta)$ है:
$A = W^2\tan(\alpha + \beta)/2$। मान लीजिए $\mathrm dA/\mathrm dt = F_a$ और $A =
0$ जब $t = 0$, तो $A = F_at$, जिससे $W$ प्राप्त होता है। $\tan\ang{6} = 0.105$ लेने पर:
$W = \sqrt{2\times40\times2/0.105} = \qty{39}{km}$, \qty{2}{Myr} के बाद; $t =
60^2\times0.105/(2\times40) = \qty{4.7}{Myr}$, \qty{60}{km} तक पहुँचने में।
\end{solution}

\begin{solution}{pb:b3:orogenic-wedges:1}
\textbf{1.} $4.5 - 1.5 = \qty{3.0}{km}$ ऊँचा; $3.0/\tan\ang{25} =
\qty{6.43}{km}$ क्षैतिज लंबाई; $3.0/\sin\ang{25} = \qty{7.10}{km}$ नति
के अनुदिश।

\textbf{2.} \omterm{def:b3:orogenic-wedges:thrust-geometry}{रैंप} की ऊँचाई से, \qty{3.0}{km}।

\textbf{3.} $S = A/r = 16.5/3.0 = \qty{5.5}{km}$।

\textbf{4.} वही, \qty{16.5}{km^2}: \omterm{def:b3:orogenic-wedges:thrust-folds}{भ्रंश-मोड़ वलन} में अतिरिक्त क्षेत्रफल
ऊपरी \omterm{def:b3:orogenic-wedges:thrust-geometry}{समतल खंड} के ऊपर हर संस्तर का $Sr$ होता है। \omterm{def:b3:orogenic-wedges:thrust-folds}{विलगन वलन} के रूप में, $16.5/4.0 =
\qty{4.1}{km}$: एक चौथाई कम।

\textbf{5.} $4\times(6.0 - 2.5) = \qty{14.0}{km}$, जो अग्र \omterm{def:b3:orogenic-wedges:thrust-geometry}{रैंप} पर
\qty{5.5}{km} के विस्थापन का ढाई गुना है।

\textbf{6.} अग्र हॉर्स, यानी अग्रभूमि के सबसे निकट वाला: हर नया \omterm{def:b3:orogenic-wedges:thrust-geometry}{रैंप}
पिछले से आगे निचली भित्ति में बनता है। हर हॉर्स, अपने \omterm{def:b3:orogenic-wedges:thrust-geometry}{रैंप} पर
ऊपर ले जाया जाकर छत-अधिक्षेप और उसके ऊपर के आवरण को, ठीक जैसे \omterm{def:b3:orogenic-wedges:thrust-folds}{भ्रंश-मोड़ वलन}
अपनी ऊपरी भित्ति को उठाता है, इसलिए छत स्टैकिंग के ऊपर मोड़ी जाती है।

\textbf{7.} $\alpha = \arctan(2.1/60) = \ang{2.00}$; $\beta =
\arctan(4.2/60) = \ang{4.00}$; टेपर \ang{6.0}।

\textbf{8.} $2.4 + 8.2 = \qty{10.6}{km}$, जो लगभग \qty{15}{km} से कम है: आधार
अभी भी घर्षणी है, और कूलॉम सिद्धांत लागू होता है।

\textbf{9.} $\phi = \ang{45.8}$, $\sin\phi = 0.717$, $K = 5.08$। शुष्क:
$\alpha + \beta = (0.0699 + 0.85)/6.08 = \qty{0.151}{rad} = \ang{8.7}$,
$\alpha = \ang{4.7}$।

\textbf{10.} $\lambda = 1000/2500 = 0.40$: $(0.0699 + 0.51)/4.05 =
\qty{0.143}{rad} = \ang{8.2}$, $\alpha = \ang{4.2}$। प्रेक्षित \ang{2.0}
पर्वतजनक कील को प्रबल रूप से उप-क्रांतिक बना देगा, जो आगे बढ़ने के बजाय
मोटी होती जाएगी: तरल अतिदाबित होने चाहिए।

\textbf{11.} $T = \qty{0.1049}{rad}$ और $u = 1 - \lambda$ लेते हुए, $T(1 + Ku) =
\beta + \mu_bu$ से $u = (\beta - T)/(TK - \mu_b) = (0.0699 -
0.1049)/(0.5327 - 0.85) = 0.110$ प्राप्त होता है: $\lambda = 0.89$।

\textbf{12.} $\sigma_v = 2500\times9.81\times10\,600 = \qty{260}{MPa}$; $P =
0.89\times260 = \qty{231}{MPa}$; द्रवस्थैतिक $1000\times9.81\times10\,600 =
\qty{104}{MPa}$। तरल पदार्थ पानी के एक स्तंभ की तुलना में
\qty{127}{MPa} अधिक दाब वहन करते हैं।

\textbf{13.} $(0.0699 + 0.85\times0.3)/(1 + 5.08\times0.3) = 0.325/2.52 =
\qty{0.129}{rad} = \ang{7.4}$, $\alpha = \ang{3.4}$। \ang{2.0} पर कील
अब उप-क्रांतिक है: वह आगे बढ़ना रोक देती है और अपने अग्र भाग के पीछे
\omterm{def:b3:orogenic-wedges:thrust-geometry}{अनुक्रम-बाह्य अधिक्षेप} से मोटी होती जाती है जब तक उसकी सतह की ढाल \ang{3.4} तक नहीं पहुँच जाती।

\textbf{14.} $F_a = \qty{10}{km/Myr}\times\qty{4.0}{km} = \qty{40}{km^2/Myr}$,
अर्थात् पट्टी के प्रति किलोमीटर \qty{4.0e4}{m^3/yr}।

\textbf{15.} $\bar e = 40/60 = \qty{0.67}{km/Myr} = \qty{0.67}{mm/yr}$।

\textbf{16.} $40\times100 = \qty{4000}{km^3/Myr} = \qty{4.0e6}{m^3/yr}$, या
$\qty{1.0e10}{kg/yr}$ (10 मिलियन टन प्रति वर्ष)।

\textbf{17.} $\mathrm dW/\mathrm dt = F_a/(W\tan(\alpha + \beta)) =
40/(60\times0.105) = \qty{6.3}{km/Myr}$।

\textbf{18.} नहीं: स्थायी अवस्था में अपरदन को अब भी पूरा आगत फ़्लक्स
हटाना होगा, इसलिए माध्य दर \qty{0.67}{mm/yr} रहती है। परंतु पिछले भाग के नीचे
जोड़ी गई चट्टान वहाँ ऊपर उठती है: अपरदन और उत्खनन पिछले भाग पर सबसे तेज़ होंगे,
जहाँ सबसे युवा शीतलन-आयु पाई जाएँगी।

\textbf{19.} $v = \qty{0.667}{km/Myr}$: $\mathrm{Pe} = 0.667\times30/31.6 =
0.634$; $G_0 = 30\times0.634/(1 - e^{-0.634}) = 30\times1.35 =
\qty{40.5}{\celsius/km}$।

\textbf{20.} $1 - e^{-\mathrm{Pe}\,z/L} = (55/900)(1 - e^{-0.634}) = 0.0287$,
इसलिए $\mathrm{Pe}\,z/L = 0.0291$ और $z_c = \qty{1.38}{km}$; $G(z_c) =
40.5\,e^{-0.0291} = \qty{39.3}{\celsius/km}$।

\textbf{21.} $1.38/0.667 = \qty{2.07}{Myr}$, नमूने की आयु, जो
\qty{2.00}{km} पर है: माध्य ऊँचाई पर समताप रेखा का विक्षेप शून्य होता है।

\textbf{22.} $v/2\kappa = 0.667/63.1 = \qty{0.0106}{km^{-1}}$, $2\pi/\Lambda =
\qty{0.314}{km^{-1}}$: $m = 0.0106 + 0.3143 = \qty{0.325}{km^{-1}}$;
$e^{-0.325\times1.38} = 0.639$; $f = (40.5 - 6.5)\times0.639/39.3 = 0.552$।

\textbf{23.} माध्य ऊँचाई \qty{2.00}{km} और माध्य आयु
\qty{2.068}{Myr} है; $\sum\Delta t\,\Delta h = 0.823$ और $\sum\Delta t^2 =
0.552$, इसलिए ढाल \qty{1.49}{km/Myr} है: आभासी दर
\qty{1.49}{mm/yr}।

\textbf{24.} $v = 1.49\times(1 - 0.552) = \qty{0.667}{mm/yr}$।

\textbf{25.} \omterm{def:b3:orogenic-wedges:thermochronology}{ऊष्मकालमापी} और फ़्लक्स संतुलन एक-दूसरे से सहमत हैं: कील
स्थायी अवस्था में है, और आभासी ढाल को अंकित मूल्य पर लेने से दर को
दो गुने से अधिक आँका गया होता। केस्ट्रेल पर्वतजनक कील की उत्खनन-दर
\textbf{\qty[text-series-to-math=true]{0.67}{mm/yr}} है।
\end{solution}
